\documentclass[10pt]{article}
\usepackage[T1]{fontenc}
\usepackage[utf8]{inputenc}
\usepackage[french]{babel}
\usepackage{amsmath,amssymb,mathrsfs,mathpazo}
\usepackage{geometry}
\geometry{a4paper,margin=2cm}
\usepackage{graphicx,array,tabularx,tabularray,colortbl,multirow,multicol,booktabs,enumitem,pifont,ulem,fancybox,diagbox,cancel}
\usepackage{pstricks,pstricks-add,pst-plot,pst-tree,pst-node,pst-text,pst-func,pst-math,pst-bspline,pst-3dplot,pst-eucl,pst-solides3d}
\let\pscalculate\undefined
\usepackage{pst-bezier}
\usepackage{tikz,pgfplots,tkz-tab}
\usetikzlibrary{arrows,arrows.meta,patterns,calc,positioning,trees,decorations.markings,shapes.geometric,shapes.misc,shapes.arrows,intersections,decorations.pathmorphing,decorations.pathreplacing,backgrounds,shadows,babel}
\pgfplotsset{compat=1.18}
\usepgfplotslibrary{fillbetween}
\usepackage[np]{numprint}
\usepackage{listings,makecell,pas-tableur,twoopt}
\let\barre\undefined
\usepackage[pstricks]{tablvar}
\definecolor{darkgreen}{RGB}{0,100,0}
\providecommand{\red}{\color{red}}
\providecommand{\blue}{\color{blue}}
\providecommand{\green}{\color{green}}
\providecommand{\black}{\color{black}}
\providecommand{\white}{\color{white}}
\providecommand{\gray}{\color{gray}}
\providecommand{\cyan}{\color{cyan}}
\providecommand{\magenta}{\color{magenta}}
\providecommand{\yellow}{\color{yellow}}
\providecommand{\violet}{\color{violet}}
\providecommand{\orange}{\color{orange}}
\providecommand{\euro}{\texteuro}
\providecommand{\EUR}{\texteuro}
\providecommand{\ieme}{\textsuperscript{e}}
\providecommand{\degres}{\ensuremath{^\circ}}
\providecommand{\degree}{\ensuremath{^\circ}}
\providecommand{\vv}[1]{\overrightarrow{#1}}
\providecommand{\Coord}[1]{#1}
\providecommand{\decosix}{}
\providecommand{\starredbullet}{\ensuremath{\bullet}}
\pagestyle{empty}
\setlength{\parindent}{0pt}
\usepackage{adjustbox}\DeclareUnicodeCharacter{2236}{:}\DeclareUnicodeCharacter{2190}{\ensuremath{\leftarrow}}
\providecommand{\R}{\mathbb{R}}
\providecommand{\N}{\mathbb{N}}
\providecommand{\Z}{\mathbb{Z}}
\providecommand{\Q}{\mathbb{Q}}
\providecommand{\D}{\mathbb{D}}
\providecommand{\C}{\mathbb{C}}
\providecommand{\e}{{\mathrm{e}}}
\providecommand{\vect}[1]{\overrightarrow{#1}}
\providecommand{\vectt}[1]{\overrightarrow{\mathrm{#1}}}
\providecommand{\vvt}[1]{\overrightarrow{\mathrm{#1}}}
\providecommand{\np}[2][]{{#2~#1}}
\providecommand{\vv}[1]{{\overrightarrow{#1}}}
\renewcommand{\d}{{\mathrm{d}}}
\providecommand{\degres}{{^\circ}}
\providecommand{\degre}{{^\circ}}
\providecommand{\blue}{\color{blue}}
\providecommand{\red}{\color{red}}
\providecommand{\enskip}{\quad}
\providecommand{\llbracket}{\lbrack\!\lbrack}
\providecommand{\rrbracket}{\rbrack\!\rbrack}
\providecommand{\euro}{{\text{€}}}
\providecommand{\no}{{\text{n°}}}
\providecommand{\interval}[2]{{[#1;#2]}}
\providecommand{\intervFO}[2]{{[#1;#2[}}
\providecommand{\dotfill}{\ldots}

\providecommand{\ds}{\displaystyle}
\providecommand{\textding}[1]{#1}
\providecommand{\Oij}{\(\mathrm{O};\vec{\imath},\vec{\jmath}\)}
\providecommand{\Oijk}{\(\mathrm{O};\vec{\imath},\vec{\jmath},\vec{k}\)}
\providecommand{\Ouv}{\(\mathrm{O};\vec{u},\vec{v}\)}
\providecommand{\intervFF}[2]{\ensuremath{[#1;#2]}}
\renewcommand{\euro}{\ifmmode\text{\texteuro}\else\texteuro\fi}
\DeclareUnicodeCharacter{00B0}{\ifmmode^{\circ}\else\textdegree\fi}
\DeclareUnicodeCharacter{20AC}{\ifmmode\text{\texteuro}\else\texteuro\fi}
% Énoncés : Éducation nationale. Transcription et corrigés : APMEP, Denis Vergès et contributeurs. Découpage et classement : CoopMaths, Sébastien Lozano, Éric Elter et contributeurs. Adaptation de la présentation : Maths Explorer.
% CC BY-SA : https://coopmaths.fr/www/about/
% https://coopmaths.fr/alea/static/bac/2024/tex/bac_2024_11_sujet1_amsud_2_cor.tex
\begin{document}
\noindent\textit{APMEP / CoopMaths. CC BY-SA. Présentation adaptée par Maths Explorer.}\par\medskip
\medskip

On dispose de deux urnes opaques U$_1$ et U$_2$.
L'urne U$_1$ contient 4 boules noires et 6 boules blanches.
L'urne U$_2$ contient 1 boule noire et 3 boules blanches.

On considère l'expérience aléatoire suivante :
on pioche au hasard une boule dans U$_1$ que l'on place dans U$_2$, puis on pioche au hasard une boule dans U$_2$.

\begin{list}{\textbullet}{On note:}
\item $N_1$ l'évènement \og Piocher une boule noire dans l'urne U$_1$ \fg.
\item $N_2$ l'évènement \og Piocher une boule noire dans l'urne U$_2$ \fg.
\end{list}

Pour tout évènement $A$, on note $\overline{A}$ son évènement contraire.

\bigskip

\textbf{PARTIE A}

\medskip

\begin{enumerate}
\item 
	\begin{enumerate}
		\item 
Si on a pioché une boule blanche dans U$_1$ et qu'on l'a mise dans U$_2$, il y a dans l'urne U$_2$ 1 boule noire et 4 boules blanches.
La probabilité de piocher alors une boule noire dans l'urne U$_2$ est $P_{\overline{N_1}}\left (N_2\right ) = \dfrac{1}{5}=0,2$.		
		
		\item 
Dans U$_1$ il y a 4 boules noires et 6 boules blanches donc la probabilité de piocher une boule noire est $P\left (N_1\right )=\dfrac{4}{10}=0,4$.		

Si on a pioché une boule noire dans U$_1$ et qu'on l'a mise dans U$_2$, il y a dans l'urne U$_2$ 2 boules noires et 3 boules blanches.
La probabilité de piocher alors une boule noire dans l'urne U$_2$ est $P_{N_1}\left (N_2\right )= \dfrac{2}{5}=0,4$.	
		
\begin{center}
\pstree[treemode=R,nodesepA=0pt,nodesepB=2.5pt,treesep = 1cm,levelsep=2.5cm]{\TR{}}
{\pstree{\TR{$N_1$~}\ncput*{$\blue 0,4$}}
	{\TR{$N_2$}\ncput*{$\blue 0,4$}
	\TR{$\overline{N_2}$}\ncput*{$\blue 0,6$}
	}
\pstree{\TR{$\overline{N_1}$~}\ncput*{$\blue 0,6$}}
	{\TR{$N_2$}\ncput*{$0,2$}
	\TR{$\overline{N_2}$}\ncput*{$\blue 0,8$}
	}
}
\end{center}

	\end{enumerate}

\item La probabilité de piocher une boule noire dans l'urne U$_1$ et une boule noire dans l'urne U$_2$ est:
$P\left ( N_1 \cap N_2 \right ) = P\left (N_1\right ) \times P_{N_1}\left (N_2\right ) = 0,4\times 0,4 = 0,16$.

\item On cherche la probabilité de piocher une boule noire dans l'urne U$_2$.

 D'après la formule des probabilités totales:

$P\left (N_2\right ) = P\left ( N_1 \cap N_2 \right ) + P\left ( \overline{N_1} \cap N_2 \right )= 0,16 + 0,6 \times 0,2 = 0,28$.

\item On a pioché une boule noire dans l'urne U$_2$.

La probabilité d'avoir pioché une boule blanche dans l'urne U$_1$ est:

$P_{N_2}\left (\overline{N_1}\right ) = \dfrac{P\left ( N_2 \cap \overline{N_1} \right )}{P\left (N_2\right ) } = \dfrac{0,12}{0,28} = \dfrac37 \approx 0,43$.


\end{enumerate}

\bigskip

\textbf{PARTIE B}

\medskip

$n$ désigne un entier naturel non nul.

L'expérience aléatoire précédente est répétée $n$ fois de façon identique et indépendante, c'est-à-dire que les urnes U$_1$ et U$_2$ sont remises dans leur configuration initiale, avec respectivement 4 boules noires et 6 boules blanches dans l'urne U$_1$ et 1 boule noire et 3 boules blanches dans l'urne U$_2$, entre chaque expérience.

On note $X$ la variable aléatoire qui compte le nombre de fois où on pioche une boule noire dans l'urne U$_2$.

On rappelle que la probabilité de piocher une boule noire dans l'urne U$_2$ est égale à $0,28$ et celle de piocher une boule blanche dans l'urne U$_2$ est égale à $0,72$.

\medskip

\begin{enumerate}
\item 
D'après le texte, on est dans le cas de $n$ répétitions,  de façon identique et indépendante, d'une expérience qui n'a que deux issues.

Donc la variable $X$ qui donne le nombre de succès, c'est-à-dire le nombre  de fois où on pioche une boule noire dans l'urne U$_2$, suit la loi binomiale de paramètres $n$ et $p= 0,28$.

\item 
On résout l'inéquation $1 - 0,72^n \geqslant 0,9.$

$\begin{aligned}[t]
1 - 0,72^n \geqslant 0,9
& \iff 1-0,9 \geqslant 0,72^n
\iff 0,1 \geqslant 0,7^n\\
& \iff \ln\left (0,1\right ) \geqslant \ln \left ( 0,72^n\right ) \quad \left (\text{fonction ln croissante sur }[0\,;\,+\infty[\strut\right )\\
& \iff \ln\left (0,1\right ) \geqslant n\,\ln \left ( 0,72\right )\\
& \iff \dfrac{\ln\left (0,1\right )}{\ln \left ( 0,72\right )} \leqslant n \quad \left (\ln\left (0,72\right )<0\strut\right )
\end{aligned}$

$ \dfrac{\ln\left (0,1\right )}{\ln \left ( 0,72\right )} \approx 7,01$ donc le plus petit entier naturel $n$ tel que:  $1 - 0,72^n \geqslant 0,9$ est 8.

\item
$0,72$ est la probabilité de piocher une boule blanche dans l'urne U$_2$.

$0,72^n$ est la probabilité de piocher $n$ boules blanches dans l'urne U$_2$ lors des $n$ tirages.

$1-0,72^n$ est la probabilité de l'événement contraire, c'est-à-dire la probabilité de piocher au moins une boule noire lors des $n$ tirages.

Il faut donc au moins 8 tirages pour que la probabilité de piocher au moins une boule noire soit supérieure à $0,9$.

\end{enumerate}

\bigskip

\textbf{PARTIE C}

\medskip

Dans cette partie les urnes U$_1$ et U$_2$ sont remises dans leur configuration initiale, avec respectivement 4 boules noires et 6 boules blanches dans l'urne U$_1$ et 1 boule noire et 3 boules blanches dans l'urne U$_2$.

\smallskip

On considère la nouvelle expérience aléatoire suivante: on pioche simultanément deux boules dans l'urne U$_1$ que l'on place dans l'urne U$_2$, puis on pioche au hasard une boule dans l'urne U$_2$.

\medskip

\begin{enumerate}
\item
Dans l'urne U$_1$ il y a 10 boules et on en prend 2 simultanément; il y a donc $\ds\binom{10}{2}=45$ tirages possibles.

\item
Il y a 4 boules noires et 6 boules blanches dans l'urne U$_1$; le nombre de tirages possibles de deux boules simultanément dans l'urne U$_1$ contenant exactement une boule blanche et une boule noire est donc $\ds\binom{4}{1}\times \binom{6}{1} = 4\times 6 = 24$.

\item 
\begin{list}{\textbullet}{On examine les différents cas.}
\item On a pioché 2 boules noires dans U$_1$ (événement noté $NN$).

\begin{list}{$\vartriangleright$}{}
\item $P\left (NN\right ) = \dfrac{\binom{4}{2}}{45} = \dfrac{6}{45}=\dfrac{2}{15}$
\item On a mis les 2 boules noires piochées dans U$_1$ dans l'urne U$_2$; l'urne U$_2$ contient alors 3 boules noires et 3 boules blanches. 

La probabilité de piocher une boule noire dans U$_2$ est alors $\dfrac{3}{6}=\dfrac{1}{2}$.
\end{list}

\item On a pioché 1 boule noire et 1 boule blanche dans U$_1$ (événement noté $NB$).

\begin{list}{$\vartriangleright$}{}
\item $P\left (NB\right ) = \dfrac{24}{45} =\dfrac{8}{15}$
\item On a mis les 2 boules,  une noire une blanche, piochées dans U$_1$ dans l'urne U$_2$; l'urne U$_2$ contient alors 2 boules noires et 4 boules blanches. 

La probabilité de piocher une boule noire dans U$_2$ est alors $\dfrac{2}{6}=\dfrac{1}{3}$.
\end{list}

\item On a pioché 2 boules blanches dans U$_1$ (événement noté $BB$).

\begin{list}{$\vartriangleright$}{}
\item $P\left (BB\right ) = \dfrac{\binom{6}{2}}{45} = \dfrac{15}{45}=\dfrac{1}{3}$
\item On a mis les 2 boules blanches piochées dans U$_1$ dans l'urne U$_2$; l'urne U$_2$ contient alors 1 boule noire et 5 boules blanches. 

La probabilité de piocher une boule noire dans U$_2$ est alors $\dfrac{1}{6}$.
\end{list}
\end{list}


On résume la situation dans un arbre pondéré.

\begin{center}
\medskip
\psset{levelsep=3cm,nodesepB=4pt, treesep=10mm}
\pstree[treemode=R,nodesepA=0pt]
       {\TR{}}
       {
       \pstree[nodesepA=4pt]{\TR{$NN$}\ncput*{$\frac{2}{15}$}}
	                        {
	                        \TR{$N_2$}~{\blue $\blacktriangleright \frac{2}{15}\times\frac{1}{2}=\frac{1}{15}=\frac{6}{90}$}\ncput*{$\frac{1}{2}$}
			                \TR{$\overline{N_2}$}\ncput*{$\frac{1}{2}$}
	                        }
       \pstree[nodesepA=4pt]{\TR{$NB$}\ncput*{$\frac{8}{15}$}}
	                        {
	                        \TR{$N_2$}~{\blue $\blacktriangleright \frac{8}{15}\times\frac{1}{3}=\frac{8}{45}=\frac{16}{90}$}\ncput*{$\frac{1}{3}$}
			                \TR{$\overline{N_2}$}\ncput*{$\frac{2}{3}$}
	                        }
       \pstree[nodesepA=4pt]{\TR{$BB$}\ncput*{$\frac{1}{3}$}}
	                        {
	                        \TR{$N_2$}~{\blue $\blacktriangleright \frac{1}{3}\times\frac{1}{6}=\frac{1}{18}=\frac{5}{90}$}\ncput*{$\frac{1}{6}$}
			                \TR{$\overline{N_2}$}\ncput*{$\frac{5}{6}$}
	                        }	                        
      }
\medskip
\end{center}

La probabilité de piocher une boule noire dans U$_2$ est, d'après la formule des probabilités totales:

$P\left (N_2\right )
= P\left ( NN \cap N_2\right ) + P\left ( NB \cap N_2\right )  + P\left ( BB \cap N_2\right ) 
= \dfrac{6}{90} + \dfrac{16}{90} + \dfrac{5}{90}
= \dfrac{27}{90} = \dfrac{3}{10} = 0,3$

La probabilité de piocher une boule noire dans l'urne U$_2$ avec cette nouvelle expérience est égale à $0,3$;  elle est donc supérieure à la probabilité de tirer une boule noire dans l'urne U$_2$ avec l'expérience de la partie A qui était de $0,28$.

\end{enumerate}

\vspace{0,5cm}

\end{document}
